% !TeX program = pdflatex
% Documentation for Overleaf/pdfLaTeX.
% This file intentionally avoids fontspec and is compatible with pdfLaTeX.
\documentclass[12pt,a4paper]{article}

\usepackage[utf8]{inputenc}
\usepackage{cmap}
\usepackage[T2A]{fontenc}
\usepackage[english,bulgarian]{babel}
\selectlanguage{bulgarian}
\usepackage{geometry}
\geometry{margin=2.5cm}
\usepackage{microtype}
\usepackage{setspace}
\onehalfspacing
\setlength{\parindent}{0pt}
\setlength{\parskip}{6pt plus 2pt}
\usepackage{xcolor}
\usepackage{graphicx}
\usepackage{tikz}
\usepackage{float}
\usepackage{caption}
\usepackage{amsmath}
\usepackage{longtable,booktabs,array,multirow}
\usepackage{etoolbox}
\AtBeginEnvironment{longtable}{\small}
\setlength{\LTleft}{0pt}
\setlength{\LTright}{0pt}
\usepackage{enumitem}
\setlist{itemsep=2pt,topsep=4pt}
\usepackage{fancyvrb}
\usepackage{fvextra}
\DefineVerbatimEnvironment{verbatim}{Verbatim}{breaklines=true,breakanywhere=true,fontsize=\small}

\providecommand{\tightlist}{\setlength{\itemsep}{0pt}\setlength{\parskip}{0pt}}
\providecommand{\pandocbounded}[1]{#1}
\providecommand{\passthrough}[1]{#1}
\sloppy
\hbadness=10000
\hfuzz=5pt

\addto\captionsbulgarian{%
  \renewcommand{\contentsname}{Съдържание}%
  \renewcommand{\figurename}{Фигура}%
  \renewcommand{\tablename}{Таблица}%
  \renewcommand{\abstractname}{Резюме}%
  \renewcommand{\refname}{Литература}%
}

\title{\textbf{Документация за курсов проект по „Бази от знания“:\\ NutriEvidence Knowledge Graph и node2vec}}
\author{
Явор Чамов\\
\small{Факултетен №: 4MI3400787}\\[0.5em]
}
\date{\today}

\begin{document}
\maketitle

\begin{abstract}
Настоящият документ описва частта от проекта \textit{NutriEvidence Agent}, която е свързана с дисциплината „Бази от знания“. Акцентът е поставен върху моделирането на биомедицински знания чрез Article--MeSH knowledge graph, представянето на статии и медицински термини като графови възли, обучението на node2vec ембединги и използването им за препоръчване на сходни научни публикации.
\end{abstract}

\tableofcontents
\newpage

\section{1. Резюме / Abstract}

\textbf{NutriEvidence Knowledge Graph} представлява модул от проекта \textit{NutriEvidence Agent}, който моделира връзките между PubMed статии и техните MeSH термини като база от знания. Вместо публикациите да се разглеждат само като независими текстови документи, те се представят като част от графова структура, в която статии, медицински понятия и тематични категории са свързани чрез експлицитни релации.

Основната идея е, че биомедицинската литература съдържа не само текстова информация, но и структурирани знания. MeSH термините са контролирани медицински понятия, чрез които публикациите се индексират в PubMed. Когато статия е свързана с даден MeSH термин, това означава, че тя разглежда конкретна медицинска тема, заболяване, популация, метод или изследователски контекст.

В проекта тези връзки се използват за изграждане на \textbf{Article--MeSH knowledge graph}. Върху графа се обучава \textbf{node2vec} модел, който генерира векторни представяния на възлите. След това embeddings на article възлите се използват за намиране на сходни статии чрез cosine similarity.

Така системата добавя графово знание към препоръчването на научни публикации. Това е особено подходящо за предмета „Бази от знания“, защото проектът демонстрира как домейн знания могат да бъдат моделирани, съхранявани, обработвани и използвани за изводи и препоръки.

\vspace{0.5em}\hrule\vspace{1.0em}

\section{2. Проблем, който проектът решава}

PubMed съдържа голямо количество биомедицински публикации. Всяка публикация може да бъде свързана с множество медицински понятия, като заболявания, симптоми, възрастови групи, методи на изследване, хранителни фактори и клинични показатели. Ако тези връзки се разглеждат само като списък от ключови думи, се губи важна структурна информация.

Основните проблеми са:

\begin{itemize}
\tightlist
\item
  трудно е да се открият статии, които са тематично близки, но не използват еднакви думи в заглавието или abstract-а;
\item
  semantic similarity върху текст може да пропусне връзки, които са видими през общи MeSH термини;
\item
  простото съвпадение на MeSH термини не отчита структурата на целия набор от статии;
\item
  липсва ясно представяне на знанията като възли, релации и домейн понятия;
\item
  потребителят има нужда от обяснима препоръка, основана не само на текст, но и на концептуална близост.
\end{itemize}

Проектът решава тези проблеми чрез изграждане на графова база от знания, в която статиите и MeSH термините са свързани. Върху този граф се прилага node2vec, за да се получи числово представяне на структурната близост между публикациите.

\vspace{0.5em}\hrule\vspace{1.0em}

\section{3. Основна цел на проекта за предмета „Бази от знания“}

Основната цел е да се разработи и документира knowledge graph модул, който:

\begin{itemize}
\tightlist
\item
  представя PubMed статии като възли в граф;
\item
  представя MeSH термините като отделни концептуални възли;
\item
  създава релации между статии и MeSH термини;
\item
  използва графовата структура като база от знания;
\item
  обучава node2vec embeddings върху графа;
\item
  намира сходни статии чрез graph-based similarity;
\item
  позволява сравнение между MeSH overlap baseline и node2vec подход;
\item
  подпомага препоръчването на научна литература чрез структурирани знания.
\end{itemize}

Фокусът не е върху LLM слоя, генерацията на отговори или медицински съвети. Фокусът е върху моделирането и използването на знания чрез граф.

\vspace{0.5em}\hrule\vspace{1.0em}

\section{4. Какво представлява базата от знания в проекта}

В контекста на проекта базата от знания е графова структура, която описва връзките между научни статии и медицински понятия. Тя може да бъде разглеждана като опростен knowledge graph, в който фактите са представени чрез тройки от вида:

\begin{verbatim}
subject --relation--> object
\end{verbatim}

За проекта основният тип тройка е:

\begin{verbatim}
Article --hasMeSHTerm--> MeSHTerm
\end{verbatim}

Пример:

\begin{verbatim}
article:37656239 --hasMeSHTerm--> mesh:cerebral_palsy
article:37656239 --hasMeSHTerm--> mesh:child
article:37656239 --hasMeSHTerm--> mesh:body_composition
\end{verbatim}

Тази структура позволява системата да отговори на въпроси като:

\begin{itemize}
\tightlist
\item
  кои статии са свързани с даден MeSH термин;
\item
  кои статии споделят общи медицински понятия;
\item
  кои MeSH термини са най-често срещани в корпуса;
\item
  кои статии са графово близки до избрана seed статия;
\item
  кои теми свързват две публикации.
\end{itemize}

\vspace{0.5em}\hrule\vspace{1.0em}

\section{5. Представяне на знанията}

\subsection{5.1. Типове възли}

Графът съдържа два основни типа възли:

\begin{longtable}{p{0.25\textwidth}p{0.65\textwidth}}
\toprule
\textbf{Тип възел} & \textbf{Описание} \\
\midrule
Article & Представя PubMed публикация. Всеки article възел се идентифицира чрез PMID. \\
MeSHTerm & Представя медицински термин от MeSH речника, свързан с една или повече статии. \\
\bottomrule
\end{longtable}

Примерни article възли:

\begin{verbatim}
article:37656239
article:12345678
article:98765432
\end{verbatim}

Примерни MeSH възли:

\begin{verbatim}
mesh:cerebral_palsy
mesh:nutritional_status
mesh:child
mesh:body_composition
\end{verbatim}

\subsection{5.2. Типове релации}

Основната релация в текущия MVP е:

\begin{longtable}{p{0.3\textwidth}p{0.6\textwidth}}
\toprule
\textbf{Релация} & \textbf{Значение} \\
\midrule
hasMeSHTerm & Свързва PubMed статия с MeSH термин, който описва съдържанието й. \\
\bottomrule
\end{longtable}

Пример:

\begin{verbatim}
article:37656239 --hasMeSHTerm--> mesh:cerebral_palsy
\end{verbatim}

В бъдеща версия могат да се добавят и други релации, например:

\begin{itemize}
\tightlist
\item
  \texttt{publishedIn} --- връзка между статия и списание;
\item
  \texttt{writtenBy} --- връзка между статия и автор;
\item
  \texttt{hasPublicationType} --- връзка между статия и тип публикация;
\item
  \texttt{broaderThan} / \texttt{narrowerThan} --- йерархични връзки между MeSH термини.
\end{itemize}

\subsection{5.3. Данни, използвани за графа}

За всяка статия се използват следните полета:

\begin{verbatim}
pmid
title
abstract
year
journal
publication_types
mesh_terms
\end{verbatim}

Най-важното поле за knowledge graph частта е:

\begin{verbatim}
mesh_terms
\end{verbatim}

То съдържа списък от медицински понятия, чрез които статията е индексирана.

\vspace{0.5em}\hrule\vspace{1.0em}

\section{6. Изграждане на Article--MeSH knowledge graph}

Процесът на изграждане на графа преминава през няколко стъпки.

\subsection{6.1. Зареждане на статии}

Системата зарежда локално кеширани PubMed статии от JSON файл:

\begin{verbatim}
data/pubmed_articles.json
\end{verbatim}

Всяка статия трябва да има поне PMID и списък от MeSH термини. Статии без MeSH термини могат да бъдат пропуснати или включени само като изолирани article възли.

\subsection{6.2. Нормализация на MeSH термините}

Преди добавяне към графа MeSH термините се нормализират, за да се избегнат дублирания. Примерни операции:

\begin{itemize}
\tightlist
\item
  премахване на излишни интервали;
\item
  превръщане към lowercase формат;
\item
  замяна на интервали със символ \texttt{\_};
\item
  премахване на празни стойности.
\end{itemize}

Пример:

\begin{verbatim}
"Cerebral Palsy" -> "cerebral_palsy"
"Body Composition" -> "body_composition"
\end{verbatim}

\subsection{6.3. Добавяне на възли и ребра}

За всяка статия се създава article възел:

\begin{verbatim}
article:<PMID>
\end{verbatim}

За всеки MeSH термин се създава MeSH възел:

\begin{verbatim}
mesh:<normalized_term>
\end{verbatim}

След това се добавя ребро между статията и всеки неин MeSH термин:

\begin{verbatim}
article:<PMID> --hasMeSHTerm--> mesh:<term>
\end{verbatim}

\subsection{6.4. Общ алгоритъм}

\begin{verbatim}
for article in articles:
    article_node = "article:" + article.pmid
    add_node(article_node, type="Article")

    for mesh_term in article.mesh_terms:
        mesh_node = "mesh:" + normalize(mesh_term)
        add_node(mesh_node, type="MeSHTerm")
        add_edge(article_node, mesh_node, relation="hasMeSHTerm")
\end{verbatim}

Резултатът е двуделен граф, в който article възлите са свързани с MeSH възли. Две статии не са директно свързани, но могат да бъдат близки чрез общи или сходни MeSH термини.

\subsection{6.5. Визуализация на Article--MeSH графа}

За демонстрация проектът включва интерактивна HTML визуализация:

\begin{verbatim}
nutri-evidence-agent/docs/article_mesh_graph.html
\end{verbatim}

Тъй като pdfLaTeX не може директно да вгражда HTML файлове, в документа е включена статична TikZ версия на същата визуализация. Сините възли са PubMed статии, зелените възли са MeSH термини, а ребрата показват релацията \texttt{hasMeSHTerm}.

\begin{figure}[H]
\centering
\input{nutri-evidence-agent/docs/fig-article-mesh-graph.tex}
\caption{Статична версия на интерактивната Article--MeSH визуализация от \texttt{docs/article\_mesh\_graph.html}.}
\label{fig:article-mesh-graph}
\end{figure}

\vspace{0.5em}\hrule\vspace{1.0em}

\section{7. MeSH overlap baseline}

Преди използване на node2vec е полезно да се реализира прост baseline, базиран на припокриване между MeSH термините на две статии.

Ако имаме две статии $A$ и $B$, техните множества от MeSH термини са $M_A$ и $M_B$. Сходството може да се изчисли чрез Jaccard similarity:

\begin{equation}
J(A, B) = \frac{|M_A \cap M_B|}{|M_A \cup M_B|}
\end{equation}

Ако две статии имат много общи MeSH термини, тяхната Jaccard стойност ще бъде по-висока.

Предимства на baseline подхода:

\begin{itemize}
\tightlist
\item
  лесен е за имплементация;
\item
  обясним е;
\item
  показва директно кои термини се припокриват;
\item
  служи като база за сравнение с node2vec.
\end{itemize}

Ограничения:

\begin{itemize}
\tightlist
\item
  отчита само директно съвпадение на термини;
\item
  не използва глобалната структура на графа;
\item
  не различава често срещани общи термини от по-специфични термини;
\item
  не намира индиректна близост между статии.
\end{itemize}

\vspace{0.5em}\hrule\vspace{1.0em}

\section{8. node2vec върху knowledge graph}

\subsection{8.1. Идея на node2vec}

\textbf{node2vec} е алгоритъм за създаване на векторни представяния на възли в граф. Основната идея е графът да се обходи чрез множество random walks. Получените последователности от възли се третират подобно на изречения, а възлите се третират подобно на думи. След това модел, подобен на Word2Vec, се обучава да поставя близки в графа възли близо един до друг във векторното пространство.

В контекста на проекта това означава:

\begin{itemize}
\tightlist
\item
  статии, които споделят подобни MeSH термини, получават близки embeddings;
\item
  статии, които са свързани чрез общи тематични области, могат да бъдат близки дори без директно припокриване на всички термини;
\item
  графовата структура се превръща в числово представяне, подходящо за similarity search.
\end{itemize}

\subsection{8.2. Random walks}

node2vec генерира множество random walks от всеки възел. Примерна разходка в Article--MeSH граф може да изглежда така:

\begin{verbatim}
article:37656239
mesh:cerebral_palsy
article:12345678
mesh:nutritional_status
article:98765432
\end{verbatim}

Така моделът научава, че тези статии и термини се срещат в сходен графов контекст.

\subsection{8.3. Параметри на node2vec}

Основни параметри:

\begin{longtable}{p{0.25\textwidth}p{0.65\textwidth}}
\toprule
\textbf{Параметър} & \textbf{Описание} \\
\midrule
Dimensions & Размерност на embedding вектора, например 64 или 128. \\
Walk length & Дължина на всяка random walk последователност. \\
Number of walks & Брой random walks, стартирани от всеки възел. \\
Window size & Размер на контекстния прозорец при обучението. \\
p & Return parameter, който контролира вероятността за връщане към предишния възел. \\
q & In-out parameter, който контролира дали обхождането е по-локално или по-глобално. \\
\bottomrule
\end{longtable}

Примерна конфигурация:

\begin{verbatim}
dimensions = 64
walk_length = 30
num_walks = 100
window = 10
p = 1.0
q = 1.0
\end{verbatim}

При $p=1$ и $q=1$ random walk поведението е балансирано. При други стойности може да се насърчи по-локално или по-глобално изследване на графа.

\subsection{8.4. Обучение на embeddings}

След генериране на random walks се обучава embedding модел. Всеки възел получава вектор:

\begin{verbatim}
article:37656239 -> [0.12, -0.04, 0.33, ...]
mesh:cerebral_palsy -> [0.08, 0.21, -0.10, ...]
\end{verbatim}

За препоръчване се използват само embeddings на article възлите, защото крайната цел е системата да препоръчва статии.

\subsection{8.5. Визуализация на node2vec embedding пространството}

В проекта са добавени и интерактивни HTML визуализации, които показват научените node2vec embeddings:

\begin{verbatim}
nutri-evidence-agent/docs/node2vec_embedding_overview.html
nutri-evidence-agent/docs/node2vec_seed_37656239.html
nutri-evidence-agent/docs/node2vec_seed_35867728.html
\end{verbatim}

Тези визуализации не показват директно суровите Article--MeSH ребра. Вместо това те показват проекция на научените node2vec вектори в двумерно пространство чрез PCA. Близки точки в тези графики означават, че съответните article възли са близки в графовото embedding пространство.

\begin{figure}[H]
\centering
\input{nutri-evidence-agent/docs/fig-node2vec-overview.tex}
\caption{Обща PCA проекция на article node2vec embeddings. Цветовете групират статиите според първоначалната PubMed заявка, а линиите показват висока graph embedding близост.}
\label{fig:node2vec-overview}
\end{figure}

\begin{figure}[H]
\centering
\input{nutri-evidence-agent/docs/fig-node2vec-seed-37656239.tex}
\caption{Seed-focused node2vec визуализация за PMID 37656239. Оранжевият възел е избраната seed статия, а останалите възли са нейни близки съседи в graph embedding пространството.}
\label{fig:node2vec-seed-37656239}
\end{figure}

\begin{figure}[H]
\centering
\input{nutri-evidence-agent/docs/fig-node2vec-seed-35867728.tex}
\caption{Seed-focused node2vec визуализация за PMID 35867728. Графиката показва как node2vec групира статии за cerebral palsy, nutrition и близки MeSH контексти.}
\label{fig:node2vec-seed-35867728}
\end{figure}

\vspace{0.5em}\hrule\vspace{1.0em}

\section{9. Graph-based recommender}

Graph-based recommender използва node2vec embeddings, за да намира статии, които са близки в графовото пространство.

\subsection{9.1. Recommend by selected article}

В този режим потребителят избира seed статия. Системата намира нейния article node:

\begin{verbatim}
article:<seed_pmid>
\end{verbatim}

След това изчислява cosine similarity между embedding-а на seed статията и embeddings на всички останали article възли.

\begin{equation}
\operatorname{sim}(A, B) = \frac{\vec{v}_A \cdot \vec{v}_B}{\|\vec{v}_A\| \|\vec{v}_B\|}
\end{equation}

Крайните резултати се сортират по similarity score и се връщат Top K най-близки статии.

\subsection{9.2. Общ процес}

\begin{verbatim}
Selected seed article
|
Find article node in graph
|
Load node2vec embedding for seed node
|
Compare with all article embeddings
|
Sort by cosine similarity
|
Return Top K similar articles
\end{verbatim}

\subsection{9.3. Изход от recommender-а}

Всеки резултат съдържа:

\begin{verbatim}
pmid
title
year
journal
mesh_terms
graph_score
method
shared_mesh_terms
\end{verbatim}

Полето \texttt{shared\_mesh\_terms} е полезно за обяснение на препоръката, защото показва кои медицински понятия свързват seed статията с препоръчаната статия.

\vspace{0.5em}\hrule\vspace{1.0em}

\section{10. Архитектура на knowledge graph модула}

Планираната структура на проекта за графовата част е:

\begin{verbatim}
src/
  graph/
    article_mesh_graph_builder.py
    node2vec_trainer.py
    graph_recommender.py
    graph_artifacts.py
  recommenders/
    mesh_overlap_recommender.py
    hybrid_recommender.py
  evaluation/
    graph_evaluator.py
\end{verbatim}

\subsection{10.1. ArticleMeshGraphBuilder}

Този компонент отговаря за:

\begin{itemize}
\tightlist
\item
  зареждане на статии;
\item
  нормализиране на MeSH термините;
\item
  създаване на article и MeSH възли;
\item
  добавяне на релации \texttt{hasMeSHTerm};
\item
  записване на графа като artifact.
\end{itemize}

\subsection{10.2. Node2VecTrainer}

Този компонент:

\begin{itemize}
\tightlist
\item
  приема готов NetworkX граф;
\item
  стартира node2vec обучение;
\item
  създава embeddings за всички възли;
\item
  запазва модела и embedding файловете;
\item
  позволява повторно използване без ново обучение.
\end{itemize}

\subsection{10.3. GraphRecommender}

Този компонент:

\begin{itemize}
\tightlist
\item
  зарежда node2vec embeddings;
\item
  филтрира само article възли;
\item
  приема seed PMID;
\item
  изчислява cosine similarity;
\item
  връща Top K graph-based препоръки.
\end{itemize}

\vspace{0.5em}\hrule\vspace{1.0em}

\section{11. Хипотези}

\subsection{11.1. Хипотеза H1: MeSH термините са подходяща основа за база от знания}

\textbf{Твърдение:} PubMed статии и техните MeSH термини могат да бъдат представени като графова база от знания, която отразява домейн връзки в биомедицинската литература.

\textbf{Очакван резултат:} Статии от близки теми ще бъдат свързани чрез общи MeSH възли или чрез близки пътища в графа.

\textbf{Критерий за отхвърляне:} Ако графът е прекалено разреден, съдържа твърде много изолирани възли или не създава полезни връзки между статии, хипотезата се отслабва.

\subsection{11.2. Хипотеза H2: node2vec улавя структурна близост между статии}

\textbf{Твърдение:} node2vec embeddings върху Article--MeSH граф могат да откриват тематично сходни статии чрез структурата на графа.

\textbf{Очакван резултат:} Graph-based recommender ще връща статии със сходни медицински теми, дори когато текстовете им не са напълно еднакви.

\textbf{Критерий за отхвърляне:} Ако node2vec резултатите са по-слаби от простия MeSH overlap baseline при повечето тестови случаи, хипотезата се отхвърля за текущия dataset.

\subsection{11.3. Хипотеза H3: Графовият подход е обясним}

\textbf{Твърдение:} Препоръките от графовия recommender могат да бъдат обяснени чрез общи MeSH термини и пътища в графа.

\textbf{Очакван резултат:} За всяка препоръчана статия системата може да покаже кои MeSH термини я свързват със seed статията.

\textbf{Критерий за отхвърляне:} Ако node2vec връща сходни статии без видими общи понятия или обясними връзки, обяснимостта на подхода е ограничена.

\vspace{0.5em}\hrule\vspace{1.0em}

\section{12. Оценяване}

Оценяването на knowledge graph частта може да се извърши чрез сравнение на няколко метода:

\begin{itemize}
\tightlist
\item
  MeSH overlap baseline;
\item
  Graph node2vec recommender;
\item
  Hybrid recommender, който комбинира semantic и graph score.
\end{itemize}

\subsection{12.1. Evaluation scenario}

Основният сценарий за оценка е \textbf{Recommend by selected article}. Потребителят избира seed article, а системата трябва да намери сходни статии.

Пример:

\begin{verbatim}
Seed article:
PMID 37656239
Low skeletal muscle mass and liver fibrosis in children with cerebral palsy.
\end{verbatim}

Системата сравнява резултатите от различните методи и проверява кои от тях връщат най-релевантни публикации.

\subsection{12.2. Метрики}

Използват се ranking метрики:

\begin{verbatim}
Precision@5
Precision@10
nDCG@5
nDCG@10
MRR
\end{verbatim}

Релевантността може да бъде зададена по скала:

\begin{verbatim}
0 = not relevant
1 = somewhat relevant
2 = relevant
3 = highly relevant
\end{verbatim}

За binary метрики може да се използва праг:

\begin{verbatim}
relevance >= 2 = relevant
\end{verbatim}

\subsection{12.3. Примерна таблица за резултати}

\begin{longtable}{p{0.34\textwidth}rrrrr}
\toprule
\textbf{Метод} & \textbf{P@5} & \textbf{P@10} & \textbf{nDCG@5} & \textbf{nDCG@10} & \textbf{MRR} \\
\midrule
MeSH overlap baseline & 0.200 & 0.200 & 0.678 & 0.678 & 0.200 \\
Graph node2vec & 0.200 & 0.200 & 0.821 & 0.775 & 0.500 \\
Hybrid & 0.200 & 0.200 & 0.770 & 0.763 & 0.250 \\
\bottomrule
\end{longtable}

Таблицата показва примерен формат за представяне на резултати. При финалната версия стойностите трябва да се изчислят върху реалните evaluation файлове на проекта.

\vspace{0.5em}\hrule\vspace{1.0em}

\section{13. Реализация}

\subsection{13.1. Използвани технологии}

За knowledge graph частта се използват:

\begin{verbatim}
Python 3.10+
pandas
numpy
NetworkX
node2vec
gensim
scikit-learn
\end{verbatim}

NetworkX се използва за построяване и съхранение на графа. node2vec и gensim се използват за обучение на embeddings. scikit-learn се използва за cosine similarity и ranking.

\subsection{13.2. Данни и artifacts}

Графовата част може да генерира следните artifacts:

\begin{verbatim}
data/artifacts/article_mesh_graph.gpickle
data/artifacts/node2vec_model.model
data/artifacts/article_node_embeddings.npy
data/artifacts/article_node_mapping.json
\end{verbatim}

Тези файлове позволяват графът и embeddings да се използват повторно без всяко стартиране да обучава модела отначало.

\subsection{13.3. Псевдокод на GraphRecommender}

\begin{verbatim}
def recommend_similar_articles(seed_pmid, top_k):
    seed_node = "article:" + seed_pmid

    if seed_node not in embeddings:
        return []

    seed_vector = embeddings[seed_node]
    scores = []

    for node, vector in embeddings.items():
        if not node.startswith("article:"):
            continue
        if node == seed_node:
            continue

        score = cosine_similarity(seed_vector, vector)
        scores.append((node, score))

    scores.sort(key=lambda x: x[1], reverse=True)
    return scores[:top_k]
\end{verbatim}

\vspace{0.5em}\hrule\vspace{1.0em}

\section{14. Ограничения}

Knowledge graph подходът има следните ограничения:

\begin{itemize}
\tightlist
\item
  качеството зависи от наличието и пълнотата на MeSH термините;
\item
  нови PubMed статии може да нямат пълни MeSH annotations;
\item
  Article--MeSH графът е сравнително прост и не използва пълната MeSH йерархия;
\item
  node2vec може да намира структурна близост, която не винаги означава реална научна релевантност;
\item
  често срещани MeSH термини като \texttt{Humans} или \texttt{Child} могат да създадат твърде общи връзки;
\item
  random walk параметрите влияят върху качеството на embeddings;
\item
  графът не оценява методологично качество, risk of bias или сила на доказателствата.
\end{itemize}

\vspace{0.5em}\hrule\vspace{1.0em}

\section{15. Бъдеща работа}

В бъдеща версия knowledge graph частта може да бъде разширена чрез:

\begin{itemize}
\tightlist
\item
  добавяне на MeSH hierarchy relations;
\item
  използване на RDF triples;
\item
  SPARQL заявки върху графа;
\item
  добавяне на authors, journals и publication types като възли;
\item
  претегляне на ръбовете според важността на MeSH термините;
\item
  филтриране на твърде общи MeSH terms;
\item
  сравнение между node2vec, DeepWalk и GraphSAGE;
\item
  визуализация на графа;
\item
  генериране на обяснения чрез shortest paths между статии;
\item
  комбиниране на graph embeddings със semantic embeddings.
\end{itemize}

\vspace{0.5em}\hrule\vspace{1.0em}

\section{16. Очакван принос на проекта}

Проектът демонстрира как знания от биомедицинска област могат да бъдат представени и използвани чрез графова база от знания. Основният принос е в комбинирането на:

\begin{verbatim}
Knowledge representation
Article--MeSH graph
Controlled medical vocabulary
Graph-based recommendation
node2vec embeddings
Similarity search
Explainable recommendations
Evaluation with ranking metrics
\end{verbatim}

Това прави проекта подходящ за дисциплината „Бази от знания“, защото показва практически пример за изграждане, обработка и използване на структурирани знания.

\vspace{0.5em}\hrule\vspace{1.0em}

\section{17. Заключение}

Графовата част на \textit{NutriEvidence Agent} представя PubMed статии и MeSH термини като база от знания. Чрез Article--MeSH knowledge graph системата може да моделира връзки между публикации и медицински понятия. node2vec превръща тази структура във векторно пространство, което позволява намиране на сходни статии чрез graph-based similarity.

За предмета „Бази от знания“ проектът е ценен, защото демонстрира целия процес: извличане на структурирани знания, моделиране като граф, обучение на графови embeddings, препоръчване на обекти и оценяване на резултатите. Въпреки ограниченията, този подход дава основа за по-богати knowledge graph решения, включително RDF, SPARQL, йерархии на понятия и обясними пътища между публикации.

\end{document}
